\section{The nonexceptional components in rank three}

Let $O$ be an original coefficient DVR with residue characteristic
$\ell$. Fix $F.\mathrm{Ok}(C_0)$ and a rank-three space $S$ on the
given frame. Write $A=S.A$, $B=S.B$ and $\pi(\xi)=\xi\cap B$.
The three entry weights are $(1,0,-1)$, so an entry in position
$(i,j)$ has degree $w_i-w_j$. The ideals $T_+,T_-$ are generated
by the lower and upper entries, respectively. Let $E$ denote the
ideal generated by the entries of odd degree. A prime is
\emph{off triangular} if it contains neither $T_+$ nor $T_-$.
Among these primes, type I means $E$ is not contained in the
prime, and type E means $E$ is contained in it.
For an ideal $J$ write $D(J)=\operatorname{Spec}A\setminus V(J)$.
For an original local chart $L$ at $x$, write $R=L.R$,
$f:\operatorname{Spec}R\to\operatorname{Spec}A$ and
$\pi_L(H)$ for contraction of any ideal $H$ of $R$ to $B$,
so $\pi_L(P)=\pi(f(P))$ for primes.  ALC17 identifies each of the three entry
ideals through this chart.

% BEGIN U67I1
\paragraph{U67I1. A Laurent coordinate through each type-I prime.}
Suppose $\xi\subset A$ is type I. Some odd entry $u$ survives
modulo $\xi$. Its degree $\epsilon$ is $1$ or $-1$. Since $\xi$
is off triangular, choose a surviving entry $v$ with degree
$-b\epsilon$, where $b=1$ or $2$. The product $a=u^b v$ belongs
to the actual degree-zero ring $B$ by RTC7 and avoids $\pi(\xi)$.
On the open set $D(a)$ one has an isomorphism of graded algebras
\[
                     B_a[U,U^{-1}]\longrightarrow A_a,
                \qquad U\longmapsto u,\quad \deg U=\epsilon .
\]
Indeed, localizing at the degree-zero element $a$ preserves the
direct grading and its degree-zero part is $B_a$. Both $u$ and $v$
are units there. For every homogeneous $z$ of degree $h$,
$zu^{-h\epsilon}$ is in $B_a$, giving surjectivity by decomposition
into degrees. Distinct Laurent powers belong to distinct direct
summands. A relation therefore forces each coefficient to vanish
after multiplication by its unit power, giving injectivity.
This calculation allows nilpotents and zero divisors in $A$.

Now take any prime $Z\subseteq\pi(\xi)$ of $B$. Extend $ZB_a$
to the Laurent ring and contract it to $A$. The result $\eta$ is
homogeneous and prime, since its localized quotient is
$(B/Z)_a[U,U^{-1}]$. It contracts to exactly $Z$ because
$a\notin Z$. Both $u$ and $v$ survive at $\eta$, proving that
both triangular ideals still survive and that an odd entry survives.
In particular $\eta$ is type I, and the fraction field of $A/\eta$
has transcendence degree exactly one over that of $B/Z$.
No inclusion between $\eta$ and $\xi$ is needed.
% END U67I1

% BEGIN U67I2
\paragraph{U67I2. The marked local incidence of the rank-three seed.}
Under the conditions of U67I1 there exists a maximal $x$ in the core
$\mathcal C=V(\mathfrak m_B A)\cap D(T_+)$ such that for
\emph{every} original chart $L$ there is a prime $P\subset R$ with
\[
             L.\rho.\mathrm{TI}(P),\qquad \pi_L(P)=Z,\qquad
             \dim R/P=\dim B/Z+1 .
\]
Indeed $\eta$ is homogeneous and avoids $T_+$, so ACC22 gives
such a maximal $x$ above $\eta$ before the choice of $L$.
The converse dimension comparison in ACC5, applied to $\eta$ with
its computed transcendence degree one, gives a prime $P$ contracting
to $\eta$ with the asserted dimension. Its contraction to $B$ is
$Z$, and the three entry-ideal comparisons in ALC17 give the type.
In particular for the source seed $M\subseteq Z\subseteq\pi(\xi)$,
$\dim B/Z=1+3d$, with $d=F.d$, this is a good type-I local prime
of dimension $3d+2$. Here good means that its contraction contains
the specified rank-three ideal $M$. This proves the rank-three
instance of $\mathrm{seedpt}$ even when the seed's first contraction
strictly contains $Z$.
% END U67I2

% BEGIN U67I3
\paragraph{U67I3. Strict dimension cuts and nice curves for either support ideal.}
Assume in addition $F.\mathrm{Big}(C_0)$. Fix any original chart
with $T_+\not\subset x$. For any $I\subset B$ satisfying
$\mathrm{FinLam}$, and any ideal $H\subset R$ with
$I\subseteq\pi_L(H)$ and $\dim R/H\ge3d+1$, there exists a
nice prime $p\supseteq H$ in the source sense.

Write $\varpi$ for a uniformizer of $O$, and $C_{vij}$ for the
ideal of all differences of the $i$th and $j$th prescribed
characters on $N_v$. Use the finite list of forbidden ideals
\[
            T_-R,\qquad IR+\varpi R+C_{vij}
                        \quad(v\in F.V,\ i\ne j).
\]
Their quotient dimensions are strictly below $3d$. For $T_-R$
this is REX9. Any prime over a listed character ideal either
contains $T_-R$ and satisfies that same bound, or is off triangular
because $T_+R=R$. In the latter case ACC15 gives
$\dim R/P+d_v\le3d+3$, and $d_v>5$. This bounds each minimal
prime of the listed ideal and hence its whole quotient, with the
zero quotient allowed. The uniformizer maps into the maximal ideal
of $R$ by the unramified coefficient extension on the chart.
Apply ACC17 at $N=3d$ to these ideals and $H$. It gives a prime
$p\supseteq H$ of quotient dimension one, containing $\varpi$
and avoiding all the forbidden ideals. It contains $\ell$ because
$O/(\varpi)$ has characteristic $\ell$. The triangular and
all coincident-character exclusions hold: $IR+\varpi R\subseteq p$,
so each forbidden character ideal forces a difference outside $p$.
This proves every field of $\mathrm{Nice}$. The argument is valid
for any such $I$, in particular both the given modularity ideal and
the distinct ideal used to control exceptional primes.
% END U67I3

% BEGIN U67I4
\paragraph{U67I4. Transfer across the exceptional components of one chart.}
Suppose $F.\mathrm{Big}(C_0)$, $\mathrm{FinLam}(M)$ and the source hypotheses
$P1'(S,M),P23'(S,M)$ hold. Assume for now that $I_E\subset B$
satisfies $\mathrm{FinLam}(I_E)$ and is contained in the
contraction of every type-E prime of $A$.
At a chart with $T_+\not\subset x$, a good type-I prime of
dimension at least $3d+1$ makes every type-I minimal prime of
$R$ good.

Choose the sharp ideal $I^\sharp$ of REX10 on this chart. A
type-I minimal prime $P$ of $R$ contains $I^\sharp$ by its
containment off $E$, so is minimal over $I^\sharp$. The latter
component primes have dimension at least $3d+2$. They are joined
by chains in which consecutive primes $P,Q$ have
$\dim R/(P+Q)\ge3d+1$, and each is type I or thick. Thickness
includes the original exceptional type and its nonsplit
representation. The useful distinction here is that the ideal
$I_E$ is in the coefficient ring, whereas $I^\sharp$ is in the
chart ring and defines its component graph.

Goodness passes up under inclusion by contraction, and type I
passes down by its three noncontainments. A nice prime is off
triangular, hence type I or type E. Above a good ideal with
quotient dimension at least $3d+1$, U67I3 with $I=M$ supplies a
nice prime. For the other curve required by the graph, take a thick
component prime $e$ below an ideal $H$ of this same quotient
dimension. The contraction $f(e)$ is type E by ALC17. Hence
$I_E\subseteq\pi_L(e)\subseteq\pi_L(H)$. U67I3 with $I=I_E$
gives a nice prime over $H$. It contains $ER$ because it contains
$e$, so this nice prime is itself type E. This use of $I_E$ does
not assume $e$ is good for $M$.

Finally $P1'$ sends goodness at a nice type-I prime to all primes
below it. In each exceptional link of the local graph the endpoints
are nice type-E primes and the intermediate thick component prime
is contained in both. These are exactly allowed links in $P23'$,
whose intermediate prime may be any thick prime. Thus $P23'$ sends
a good type-I prime below the initial endpoint to goodness of all
type-I primes below the last endpoint, including a chain of length
zero. These facts supply every local condition of ACP14 with threshold
$3d+1$, giving the claimed local propagation.
% END U67I4

% BEGIN U67I5
\paragraph{U67I5. Coverage of all nonexceptional primes of the original ring.}
Retain the hypotheses of U67I4 and $\mathrm{SeedIn}(S,M)$ as in
U67I2. Then $M\subseteq\pi(\xi)$ for \emph{every} type-I prime
$\xi$ of $A$. We verify the geometric steps transferring
the local argument to $X=\operatorname{Spec}A$. Here
\[
 T=D(T_+)\cap D(T_-)\cap D(E),\qquad G=V(MA),\qquad
 Z_c=\{x\in\mathcal C\mid x\text{ is maximal in }X\}.
\]
By G67T4 the core is preconnected and is contained in the closure of
its part outside $V(E)$. ACC21 gives a maximal core point in
every nonempty intersection of the core with an open and a closed
set. In particular every closed set meeting the core meets $Z_c$,
and each open neighborhood of a core point outside $V(E)$
meets $Z_c\setminus V(E)$, by taking its intersection with
$D(E)$. Thus $\mathcal C\subseteq\overline{Z_c\setminus V(E)}$.
Every component $Y$ meeting $T$ has type-I generic prime by
openness. This is a homogeneous minimal prime by ACC20. ACC22
puts a core point on $Y$.

Fix $x\in Z_c$. ALC21 supplies a chart there and ALC11 gives
flatness. By ACP16 every global component through $x$ is the
closure of the contraction of a chart minimal prime. If $Y$
meets $T$, its generic prime lies in $T$ and its chart minimal
prime $P$ is type I by ALC17. It therefore contains $I^\sharp$
and has $\dim R/P\ge3d+2$ by REX10. If $Y\subseteq G$ then
$P$ is good. U67I4 makes the type-I chart minimal prime for
every other component through $x$ meeting $T$ good as well.
The contractions and closedness of $G$ make each of these whole
components good. At $x\notin V(E)$ the sharp ideal is zero and
$ER=R$, so all chart minimal primes are type I by REX10.
Taking one gives a component through $x$ which meets $T$.

The actual seed of U67I2 gives a single $x\in Z_c$ at whose
chart there is a good type-I prime of dimension $3d+2$. A
minimal prime of $R$ below it is type I by generization and good
by U67I4. Its contraction is minimal in $A$ by flatness, hence
its closure is a whole component contained in $G$ and meeting
$T$. This establishes the last of the eight incidence conditions
of ACP15. The entry ring is Noetherian by RTC7. ACP15 now
gives $T\subseteq G$, the original $\mathrm{AllGood}(S,M)$.

In the original $\mathrm{PropagationTheoremGL3}$ the additional
given rank-two space $S_2$ on this same frame has an ideal $M_2$
satisfying $\mathrm{FinLam}$ and $\mathrm{AllGood}$. The rank-two
corner construction C67E5 supplies exactly the $I_E$
used above from those two properties. ACP18 gives the single
finite-weight ideal $M\cap I_E$ for \emph{all} off-triangular
primes. Thus the two simultaneous conclusions are
$\mathrm{AllGood}(S,M)$ and $\mathrm{OffFin}(S)$, with the
original finite-weight, $P1'$, $P23'$ and seed assumptions on $M$.
ACP11 and RTC24 can be applied once the marked arithmetic
construction supplies those premises on its own spaces and ideals.
The component theorem itself does not produce the modularity
ideals or their arithmetic seed data.
% END U67I5
